\documentclass[../../main.tex]{subfiles}
\begin{document}
\begin{oframed}
  \subfile{problem}
\end{oframed}

{\bf [解答]}

\paragraph{(1)}

\begin{figure}[htbp]
  \centering
  \tdplotsetmaincoords{70}{115}
  \begin{tikzpicture}[scale=1.35, >=stealth, every node/.style={font=\small}]

    \def\Rval{2.0} % 球の半径
    \def\dval{1.1} % 原点から平面への距離
    \pgfmathsetmacro{\rval}{sqrt(\Rval*\Rval - \dval*\dval)} % 交円の半径

    % 1. 球面 A 本体のシェーディングと輪郭（スクリーン座標）
    \coordinate (O_scr) at (0, 0);
    \shadedraw[ball color=blue!12, opacity=0.35, draw=gray!75!black, line width=1.1pt] (O_scr) circle (\Rval);

    % 2. 3D座標系内の要素
    \begin{scope}[tdplot_main_coords]
      % 円の中心 H
      \pgfmathsetmacro{\hx}{\dval / 1.732}
      \pgfmathsetmacro{\hy}{\dval / 1.732}
      \pgfmathsetmacro{\hz}{\dval / 1.732}
      \coordinate (H3) at (\hx, \hy, \hz);
      \coordinate (O3) at (0, 0, 0);

      % 座標軸 (奥側)
      \draw[dashed, gray!45] (O3) -- (1.0, 0, 0);
      \draw[dashed, gray!45] (O3) -- (0, 1.0, 0);
      \draw[dashed, gray!45] (O3) -- (0, 0, 1.0);

      % 赤道（xy 平面の大円）
      \draw[gray!55, dashed] plot[domain=0:180, samples=40] ({\Rval*cos(\x)}, {\Rval*sin(\x)}, 0);
      \draw[gray!55] plot[domain=180:360, samples=40] ({\Rval*cos(\x)}, {\Rval*sin(\x)}, 0);

      % 平面 B の基底ベクトル (単位長)
      \def\pw{2.1}
      \coordinate (P1) at ({\hx + \pw*0.7071 + \pw*0.4082}, {\hy - \pw*0.7071 + \pw*0.4082}, {\hz - \pw*0.8165});
      \coordinate (P2) at ({\hx - \pw*0.7071 + \pw*0.4082}, {\hy + \pw*0.7071 + \pw*0.4082}, {\hz - \pw*0.8165});
      \coordinate (P3) at ({\hx - \pw*0.7071 - \pw*0.4082}, {\hy + \pw*0.7071 - \pw*0.4082}, {\hz + \pw*0.8165});
      \coordinate (P4) at ({\hx + \pw*0.7071 - \pw*0.4082}, {\hy - \pw*0.7071 - \pw*0.4082}, {\hz + \pw*0.8165});

      % 平面 B の半透明塗り
      \fill[teal!20, opacity=0.45] (P1) -- (P2) -- (P3) -- (P4) -- cycle;
      \draw[teal!80!black, thick] (P1) -- (P2) -- (P3) -- (P4) -- cycle;
      \node[teal!90!black, font=\footnotesize\bfseries, above right] at (P2) {平面 $B$};

      % 交円 D
      \draw[line width=1.4pt, blue!85!black, domain=0:360, samples=60]
      plot ({ \hx + \rval*cos(\x)*0.7071 + \rval*sin(\x)*0.4082 },
        { \hy - \rval*cos(\x)*0.7071 + \rval*sin(\x)*0.4082 },
      { \hz - \rval*sin(\x)*0.8165 });
      \node[blue!85!black, font=\footnotesize\bfseries, left] at ($(H3) + (-1.1, 1.1, 0)$) {円 $D$};

      % 垂線 OH
      \draw[dashed, red!85!black, thick] (O3) -- (H3) node[midway, left=1pt, font=\footnotesize] {$l$};

      % 交円上の2点 P, Q (P: 60度, Q: 150度)
      \pgfmathsetmacro{\px}{\hx + \rval*cos(60)*0.7071 + \rval*sin(60)*0.4082}
      \pgfmathsetmacro{\py}{\hy - \rval*cos(60)*0.7071 + \rval*sin(60)*0.4082}
      \pgfmathsetmacro{\pz}{\hz - \rval*sin(60)*0.8165}
      \coordinate (Ppt) at (\px, \py, \pz);

      \pgfmathsetmacro{\qx}{\hx + \rval*cos(150)*0.7071 + \rval*sin(150)*0.4082}
      \pgfmathsetmacro{\qy}{\hy - \rval*cos(150)*0.7071 + \rval*sin(150)*0.4082}
      \pgfmathsetmacro{\qz}{\hz - \rval*sin(150)*0.8165}
      \coordinate (Qpt) at (\qx, \qy, \qz);

      \draw[line width=1.2pt, orange!85!black] (Ppt) -- (Qpt);
      \node[circle, fill=orange!85!black, inner sep=1.2pt] at (Ppt) {};
      \node[circle, fill=orange!85!black, inner sep=1.2pt] at (Qpt) {};
      \node[right=2pt, font=\footnotesize] at (Ppt) {$P$};
      \node[above=2pt, font=\footnotesize] at (Qpt) {$Q$};

      % 点 O, H
      \node[circle, fill=black, inner sep=1.2pt] at (O3) {};
      \node[below left=1pt] at (O3) {$O$};
      \node[circle, fill=red!85!black, inner sep=1.2pt] at (H3) {};
      \node[above=2pt, font=\footnotesize] at (H3) {$H$};

      % 座標軸 (手前側)
      \draw[->, thick, gray!80!black] (1.0, 0, 0) -- (2.8, 0, 0) node[below left] {$x$};
      \draw[->, thick, gray!80!black] (0, 1.0, 0) -- (0, 2.8, 0) node[right] {$y$};
      \draw[->, thick, gray!80!black] (0, 0, 1.0) -- (0, 0, 2.7) node[above] {$z$};
    \end{scope}

    % 球面 A のラベル
    \node[gray!70!black, font=\footnotesize] at (-1.5, 1.4) {球面 $A$};

  \end{tikzpicture}
  \caption{球面 $A$ と平面 $B$ の交円 $D$ の様子}
  \label{fig:sphere_plane_3d}
\end{figure}

与えられた球面$A$と平面$B$
\begin{align}
  \begin{cases}
    A: x^2+y^2+z^2=\dfrac{1}{3}(r^2+2) \\
    B: x+y+z=r
  \end{cases} \label{eq:1}
\end{align}
について考える．$A$の中心 $(0,0,0)$ と平面$B$の距離$l$は，点と平面の距離公式より
\begin{align}
  l = \dfrac{r}{\sqrt{1+1+1}} = \dfrac{\sqrt{3}}{3}r
\end{align}
だから，$D$ は半径
\begin{align}
  \sqrt{\dfrac{1}{3}(r^2+2) - \left(\dfrac{\sqrt{3}}{3}r\right)^2} = \sqrt{\dfrac{2}{3}}
\end{align}
の円である．よってこの円の中の点$P,Q$について
\begin{align}
  |PQ| \le (D\text{の直径}) = 2\sqrt{\dfrac{2}{3}}
\end{align}
となり題意は示された．\qed

\begin{figure}[htbp]
  \centering
  \begin{tikzpicture}[scale=1.4, >=stealth, every node/.style={font=\small}]
    % 中心 O, 半径 R
    \def\R{2.4}
    \def\d{1.3} % OH = l = sqrt(3)/3 * r
    \pgfmathsetmacro{\rho}{sqrt(\R*\R - \d*\d)}

    \coordinate (O) at (0, 0);
    \coordinate (H) at (\d, 0);
    \coordinate (P) at (\d, \rho);
    \coordinate (Q) at (\d, -\rho);

    % 球面の断面（円）
    \draw[line width=1.2pt, gray!70!black] (O) circle (\R);
    \node[gray!70!black, font=\footnotesize] at (-1.7, 1.7) {球面 $A$ の大円};

    % 平面 B の断面（直線）
    \draw[line width=1.1pt, teal!80!black] (\d, -2.4) -- (\d, 2.4) node[above, font=\footnotesize] {平面 $B$};

    % 交円 D の直径 PQ
    \draw[line width=1.8pt, blue!85!black] (P) -- (Q);
    \node[right=2pt, blue!85!black, font=\footnotesize] at (\d, 0.7) {交円 $D$ の直径};

    % 直角三角形 OHP
    \draw[dashed, red!85!black, thick] (O) -- (H)
    node[midway, below, font=\footnotesize] {$l = \frac{\sqrt{3}}{3}r$};
    \draw[dashed, blue!85!black, thick] (O) -- (P)
    node[midway, above left, font=\footnotesize] {$R = \sqrt{\frac{r^2+2}{3}}$};
    \draw[dashed, gray!60, thick] (H) -- (P)
    node[midway, right=2pt, font=\footnotesize] {$\sqrt{\frac{2}{3}}$};

    % 直角マーク at H
    \draw[gray!60, thin] (\d-0.18, 0) -- (\d-0.18, 0.18) -- (\d, 0.18);

    % 点
    \node[circle, fill=black, inner sep=1.3pt] at (O) {};
    \node[left=2pt] at (O) {$O$};

    \node[circle, fill=red!85!black, inner sep=1.3pt] at (H) {};
    \node[below left=2pt, font=\footnotesize] at (H) {$H$};

    \node[circle, fill=blue!85!black, inner sep=1.3pt] at (P) {};
    \node[above right=2pt, font=\footnotesize] at (P) {$P$};

    \node[circle, fill=blue!85!black, inner sep=1.3pt] at (Q) {};
    \node[below right=2pt, font=\footnotesize] at (Q) {$Q$};

    % ピタゴラスの補足説明ボックス
    \node[draw=blue!70!black, fill=blue!5, rounded corners=2pt, inner sep=4pt, font=\scriptsize, align=left]
    at (-0.8, -1.2) {
      $
      \begin{aligned}
        HP &= \sqrt{OP^2 - OH^2} \\
        &= \sqrt{\frac{r^2+2}{3} - \left(\frac{\sqrt{3}}{3}r\right)^2} = \sqrt{\frac{2}{3}}
      \end{aligned}$
    };
  \end{tikzpicture}
  \caption{$O$ から平面 $B$ へ下ろした垂線を含む断面とピタゴラスの定理による交円の半径の導出}
  \label{fig:cross_section_pythagoras}
\end{figure}

\paragraph{(2)}

\cref{eq:1}で第二式を第一式に代入して$r$を消すと
\begin{align}
  &x^2 + y^2 + z^2 = \dfrac{1}{3} \left\{ (x+y+z)^2 + 2 \right\} \\
  \therefore
  &(x-y)^2 + (y-z)^2 + (z-x)^2 = 2
\end{align}
である．$x, y, z \in \mathbb{Z}$ だから，$(x-y)^2, (y-z)^2, (z-x)^2$ のうち2つが1，のこり1つが0である．すなわち
\begin{align}
  \left\{ (x-y)^2, (y-z)^2, (z-x)^2 \right\} = \{1, 1, 0\}, \{1, 0, 1\}, \{0, 1, 1\}
\end{align}
である．対称性から $x \le y \le z$ とするとこれをみたすのは
\begin{align}
  (x, y, z) = (x, x+1, x+1), (x, x, x+1)
\end{align}
のみである．従ってこれらを入れ替えたものも考えて，$k \in \mathbb{Z}$ として
\begin{align}
  \begin{cases}
    r = 3k+2 \text{ の時 } & (x, y, z) = (k, k+1, k+1), (k+1, k, k+1), (k+1, k+1, k) \\
    r = 3k+1 \text{ の時 } & (x, y, z) = (k, k, k+1), (k, k+1, k), (k+1, k, k)
  \end{cases}
\end{align}
である．

\end{document}
